\section{Layered Architecture}


Zoo Network operates as a specialized Layer 2 atop Hanzo's base compute infrastructure. This section details the multi-layer design.

\subsection{Layer 0: Lux Consensus}

\textbf{Lux} \cite{lux2025whitepaper} provides the foundational consensus layer with:

\begin{itemize}[leftmargin=*]
\item \textbf{Wave Consensus}: DAG-based finality via repeated sub-sampled voting \cite{rocket2020snowman}. Achieves sub-second confirmation with Byzantine fault tolerance.
\item \textbf{Post-Quantum Cryptography}: CRYSTALS-Dilithium signatures resist quantum attacks.
\item \textbf{Multi-Chain Architecture}: Supports heterogeneous virtual machines (EVM, WASM, custom).
\end{itemize}

Zoo inherits security and finality guarantees from Lux's validator set.

\subsection{Layer 1: Hanzo Compute Infrastructure}

\textbf{Hanzo} \cite{hanzo2025paper} extends Lux with compute-specific capabilities:

\begin{itemize}[leftmargin=*]
\item \textbf{GPU Compute Nodes}: Distributed GPU clusters run containerized LLM inference. Nodes stake collateral and earn rewards proportional to compute provided.
\item \textbf{Proof of Compute}: Validators verify inference correctness via sampling or zero-knowledge proofs.
\item \textbf{Storage Primitives}: Integration with IPFS for content-addressable data and Arweave for permanent archival.
\item \textbf{Rust Crates}: Hanzo provides reusable libraries (consensus, mining, HTTP APIs, libp2p networking) that Zoo extends.
\end{itemize}

Hanzo is blockchain-agnostic compute infrastructure; Zoo specializes it for AI/ML workloads.

\subsection{Layer 2: Zoo AI/ML Specialization}

Zoo adds AI-specific components:

\begin{enumerate}[leftmargin=*]
\item \textbf{Smart Contracts}:
   \begin{itemize}
   \item \textit{Inference Router}: Routes queries to available GPU nodes, specifies experience library versions, records proofs.
   \item \textit{Experience Registry}: Stores Merkle roots of experience libraries, handles DAO votes for curation.
   \item \textit{Governance Contract}: Manages proposals, voting, and treasury allocation.
   \end{itemize}

\item \textbf{Off-Chain Components}:
   \begin{itemize}
   \item \textit{Semantic Context Manager}: Maintains experience embeddings, performs similarity search for retrieval.
   \item \textit{Experience Storage}: IPFS (current library), Arweave (immutable history).
   \item \textit{GPU Compute Nodes}: Load frozen base models, inject experiences into context windows, execute inference.
   \end{itemize}

\item \textbf{Rust Libraries} (extend Hanzo):
   \begin{itemize}
   \item \texttt{zoo\_mcp}: Model Context Protocol for experience management.
   \item \texttt{zoo\_embedding}: Vector embedding generation.
   \item \texttt{zoo\_job\_queue\_manager}: Training job scheduling.
   \item \texttt{zoo\_tools\_primitives}: AI-specific utilities.
   \end{itemize}
\end{enumerate}

\subsection{Data Flow}

\begin{enumerate}
\item User submits query via Inference Router contract.
\item Router selects GPU node and experience library version.
\item GPU node fetches base model (frozen weights) and experience library (IPFS).
\item Experiences are injected into context window.
\item Model generates response.
\item Response returned to user; proof recorded on-chain.
\end{enumerate}

This architecture separates concerns: Lux handles consensus, Hanzo handles compute, Zoo handles AI-specific logic.

